%----------------------------------------------------------------------
%% APP: ASYMPTOTIC
%----------------------------------------------------------------------
% !TEX root = ./Main.tex


Our goal in this appendix is to prove \cref{thm:stable,cor:stable}, which we restate below for convenience:

\stable*

\corstable*

Our proof strategy will comprise the following basic steps:

\begin{enumerate}

\item
To begin with, we will show that the squared distance
\begin{equation}
\label{eq:energy}
\energy(\policy)
	= \frac{1}{2}\norm{\policy - \eq}^{2}
\end{equation}
can be seen as a ``local Lyapunov function'' for \eqref{eq:PG} in the sense that it is locally decreasing near $\eq$, up to a series of error terms \textendash\ both zero-mean and non-zero-mean.

\item
Due to these errors, the evolution of the iterates $\curr[\energy] \defeq \energy(\curr)$ of $\energy$ over time may exhibit \emph{significant} jumps:
in particular, a single ``bad'' realization of the noise could carry $\curr$ out of the basin of attraction of $\eq$, possibly never to return.
To exclude this event, our second step will be to show that the aggregation of these errors can be controlled with probability at least $1-\conf$.
% This is the most technically challenging part of our argument and it unfolds in a series of lemmas below.

\item
Conditioned on the above, we will show that, with probability at least $1-\conf$, the iterates $\curr[\energy]$ cannot grow more than a token value.
%(depending on the basin of attraction of $\eq$).
As a result, if \eqref{eq:PG} is initialized close to $\eq$, it will remain in a neighborhood thereof for all $\run$ (again, with probability at least $1-\conf$).

\item
Thanks to this ``stochastic Lyapunov stability'' result, we employ a series of martingale limit theory arguments to extract a subsequence converging to $\eq$.

\item
Finally, we show that the increments of $\curr[\energy]$ are summable;
hence, by invoking the Gladyshev's lemma \cite[p.~49]{Pol87}, we conclude that $\curr[\energy]$ converges to some (finite) random variable $\energy_{\infty}$.
Combining this fact with the existence of a convergent subsequence, we obtain the desired conclusion that $\curr$ converges to $\eq$ with probability at least $1-\conf$.
\end{enumerate}

In the sequel, we make the above precise in a series of intermediate results.


%----------------------------------------------------------------------
%% Energy
%----------------------------------------------------------------------
\subsection{Energy inequality}

We begin by establishing a ``quasi-Lyapunov'' inequality for the iterates $\curr[\energy] = \norm{\curr - \eq}^{2}/2$ of \eqref{eq:energy}.

\begin{lemma}
\label{lem:energy}
Let $\curr[\energy] \defeq \energy(\curr)$.
Then, for all $\run = \running$, we have
\begin{equation}
\label{eq:template}
\next[\energy]
	\leq \curr[\energy]
		+ \curr[\step] \braket{\gvalue(\curr)}{\curr - \eq}
		+ \curr[\step] \curr[\snoise]
		+ \curr[\step] \curr[\sbias]
		+ \curr[\step]^{2} \curr[\second]^{2},
\end{equation}
where the error terms $\curr[\snoise]$, $\curr[\sbias]$, and $\curr[\second]$ are given by
\begin{equation}
\label{eq:errors}
\curr[\snoise]
	= \braket{\curr[\noise]}{\curr - \eq},
%	\\
	\quad
\curr[\sbias]
	= \poldiam \curr[\bbound]
%	= \curr[\bbound] \norm{\\curr - \eq}
	\quad
	\text{and}
	\quad
%	\\
\curr[\second]^{2}
	= \tfrac{1}{2} \dnorm{\curr[\signal]}^{2}.
\end{equation}
with $\poldiam \defeq \max_{\policy,\policyalt\in\policies} \norm{\policy - \policyalt}$.
\end{lemma}

\begin{proof}
By the definition of the iterates of \eqref{eq:PG}, we have
\begin{align}
\next[\energy]
	= \frac{1}{2} \norm{\next - \eq}^{2}
	&= \frac{1}{2} \norm{\proj_{\policies}(\curr + \curr[\step]\curr[\signal]) - \proj_{\policies}(\eq)}^{2}
	\notag\\
	&\leq \frac{1}{2} \norm{\curr + \curr[\step]\curr[\signal] - \eq}^{2}
	\notag\\
	&= \frac{1}{2} \norm{\curr - \eq}^{2}
		+ \curr[\step] \braket{\curr[\signal]}{\curr - \eq}
		+ \frac{1}{2} \curr[\step]^{2} \norm{\curr[\signal]}^{2}
	\notag\\
	&= \curr[\energy]
		+ \curr[\step] \braket{\gvalue(\curr) + \curr[\noise] + \curr[\bias]}{\curr - \eq}
		+ \frac{1}{2} \curr[\step]^{2} \norm{\curr[\signal]}^{2}
	\notag\\
%	&= \curr[\energy] + \curr[\step]\braket{\gvalue(\curr)}{\curr - \eq} + \curr[\step]\curr[\snoise] + \curr[\step]\braket{\curr[\bias]}{\curr - \eq} + \curr[\step]^{2} \curr[\second]^{2}
%	\notag\\
	&\leq \curr[\energy]
		+ \curr[\step] \braket{\gvalue(\curr)}{\curr - \eq}
		+ \curr[\step] \curr[\snoise]
		+ \curr[\step] \curr[\sbias]
		+ \curr[\step]^{2} \curr[\second]^{2}
\end{align}
where we used the Cauchy-Schwarz inequality to bound the bias term as $\braket{\curr[\bias]}{\curr -\eq} \leq \dnorm{\curr[\bias]} \cdot \norm{\curr - \eq} \leq \poldiam \curr[\bbound] = \curr[\sbias]$.
\end{proof}


%----------------------------------------------------------------------
%% Error control
%----------------------------------------------------------------------
\subsection{Error control and stability}

The second major step in our proof (and the most challenging one from a technical standpoint) is to establish a suitable measure of control over the error increments in \eqref{eq:energy}, with the aim of showing that the process $\curr$ never leaves a neighborhood of $\eq$.

To make this idea precise, let $\basin = \setdef{\policy\in\policies}{\norm{\policy - \eq} \leq \radius}$ be a ball of radius $\radius$ based on $\eq$ in $\policies$ so that $\braket{\gvalue(\policy)}{\policy - \eq} < 0$ for all $\policy\in\basin\exclude{\eq}$ (without loss of generality, we can assume that $\basin$ is maximal in that regard).
We will then examine the event that the aggregation of the error terms in \eqref{eq:energy} is not sufficient to drive $\curr$ to escape from $\basin$.

To that end, we will begin by aggregating the errors in \eqref{eq:energy} as
\begin{equation}
\label{eq:err-agg}
\curr[\mart]
	= \sum_{\runalt=\start}^{\run} \iter[\step] \iter[\snoise]
	\quad
	\text{and}
	\quad
\curr[\submart]
	= \sum_{\runalt=\start}^{\run} \bracks{\iter[\step] \iter[\sbias] + \iter[\step]^{2} \iter[\second]^{2}}.
\end{equation}
Since $\exof{\curr[\snoise] \given \curr[\filter]} = 0$, we have $\exof{\curr[\mart] \given \curr[\filter]} = \prev[\mart]$, so $\curr[\mart]$ is a martingale;
likewise, $\exof{\curr[\submart] \given \curr[\filter]} \geq \prev[\submart]$, so $\curr[\submart]$ is a submartingale.
Then, using a technique of \citet{HIMM19} that builds on an earlier idea by \citet{MZ19}, we will also consider the ``mean square'' error process
\begin{align}
\label{eq:noise-tot}
\curr[\toterr]
	&= \curr[\mart]^{2}
		+ \curr[\submart],
\end{align}
and the associated indicator events
\begin{subequations}
\label{eq:events}
\begin{align}
%\label{eq:evt-stay}
\curr[\event]
%	\equiv \curr[\event](\basin)
	&= \braces*{\iter\in\basin \; \text{for all $\runalt=\running,\run$}}
	\quad
	\text{and}
	\quad
\curr[\eventalt]
%	\equiv \curr[\eventalt](\thres)
	= \braces*{\iter[\toterr] \leq \thres \; \text{for all $\runalt = \running,\run$}},
\end{align}
\end{subequations}
where, with a fair amount of hindsight,
the error tolerance level $\thres>0$ is such that $2\thres + \sqrt{\thres} < \radius$, and we are employing the convention $\event_{0} = \eventalt_{0} = \samples$ (since every statement is true for the elements of the empty set).
We will then assume that $\init$ is initialized in a ball of radius $\sqrt{2\thres}$ centered at $\eq$, viz.
\begin{equation}
\label{eq:nhd-init}
\nhd
	= \setdef{\policy\in\policies}{\energy(\policy) \leq \thres}
	= \setdef{\policy\in\policies}{\norm{\policy-\eq}^{2}/2 \leq \thres}.
\end{equation}
With all this in hand, the key to showing that $\curr$ remains close to $\eq$ with high probability is the following conditional estimate:

\begin{lemma}
\label{lem:events}
Let $\curr$ be the sequence of play generated by \eqref{eq:PG} initialized at $\init\in\nhd$.
We then have:
%with step-size and policy gradient estimates as per \cref{thm:stable}.
%and with step-size and policy gradient estimates such that $\pexp + \bexp >1$ and $\pexp - \sexp > 1/2$ as per \cref{thm:stable}.
%We then have:
\begin{enumerate}
\item
$\next[\event] \subseteq \curr[\event]$
and
$\next[\eventalt] \subseteq \curr[\eventalt]$
for all $\run=\running$
\item
$\prev[\eventalt] \subseteq \curr[\event]$
for all $\run=\running$
\item
Consider the ``bad realization'' event
\begin{align}
\label{eq:evt-bad}
\curr[\tilde\eventalt]
	\defeq \prev[\eventalt] \setminus \curr[\eventalt]
%	= \prev[\eventalt] \cap \{\curr[\toterr] > \thres\}
%	\notag\\
	= \braces*{\text{$\iter[\toterr] \leq \thres$ for $\runalt=\running,\run-1$ and $\curr[\toterr] > \thres$}},
\end{align}
and
let $\curr[\tilde\toterr] = \curr[\toterr] \one_{\prev[\eventalt]}$ be the cumulative error subject to the noise being ``small''.
%at stage $\run$ conditioned on the noise being ``small'' until that time.
%denote the event that $\iter\in\basin$ for all $\runalt=\running,\run-1$ but $\curr$ escapes $\basin$.
Then we have:
\begin{equation}
\label{eq:noise-tot-cond}
\exof{\curr[\tilde\toterr]}
	\leq \exof{\prev[\tilde\toterr]}
		+ \curr[\step] \poldiam \curr[\bbound]
		+ \curr[\step]^{2} \poldiam^{2} \curr[\sdev]^{2}
		+ \tfrac{3}{2}  \curr[\step]^{2} (\vbound^{2} + \curr[\bbound]^{2} + \curr[\sdev]^{2})
		- \thres \probof{\prev[\tilde\eventalt]},
\end{equation}
where, by convention,
%$\poldiam = \sup_{\dpoint} \norm{\nabla\energy(\policy)}$
$\tilde\eventalt_{0} = \varnothing$ and $\beforeinit[\tilde\toterr] = 0$.
\end{enumerate}
\end{lemma}

\begin{remark*}
In the above (and what follows), the notation $\one_{A}$ is used to indicate the logical indicator of an event $A\subseteq\samples$, \ie $\one_{A}(\sample) = 1$ if $\sample\in A$ and $\one_{A}(\sample) = 0$ otherwise.
\end{remark*}

The proof of \cref{lem:events} is quite technical, so we first proceed to derive an important stability result based on this estimate.

\begin{proposition}
\label{prop:contain}
Fix some confidence threshold $\conf>0$ and let $\curr$ be the sequence of play generated by \eqref{eq:PG} with step-size and policy gradient estimates as per \cref{thm:stable}.
%that $\pexp + \bexp >1$ and $\pexp - \sexp > 1/2$ as per \cref{thm:stable}.
We then have:
\begin{equation}
\label{eq:contain}
\probof{\curr[\eventalt] \given \init\in\nhd}
	\geq 1-\conf
	\quad
	\text{for all $\run=\running$}
\end{equation}
provided that $\step$ is small enough \textpar{or $\runoff$ large enough} relative to $\conf$.
\end{proposition}

\begin{proof}
We begin by bounding the probability of the ``bad realization'' event $\curr[\tilde\eventalt] = \prev[\eventalt] \setminus \curr[\eventalt]$.
Indeed, if $\init\in\nhd$, we have:
\begin{align}
\label{eq:prob-large1}
\probof{\curr[\tilde\eventalt]}
	&= \probof{\prev[\eventalt] \setminus \curr[\eventalt]}
%	= \probof{\prev[\eventalt] \cap \{\curr[\toterr] > \thres\}}
%	\notag\\
	= \exof{\one_{\prev[\eventalt]} \times \oneof{\curr[\toterr] > \thres}}
%	\notag\\
	\leq \exof{\one_{\prev[\eventalt]} \times (\curr[\toterr] / \thres)}
%	\notag\\
	= \exof{\curr[\tilde\toterr]} / \thres
\end{align}
where, in the penultimate step, we used the fact that $\curr[\toterr] \geq 0$ (so $\oneof{\curr[\toterr]>\thres} \leq \curr[\toterr]/\thres$).
Telescoping \eqref{eq:noise-tot-cond} then yields
\begin{equation}
\label{eq:prob-large2}
\exof{\curr[\tilde\toterr]}
	\leq \exof{\beforeinit[\tilde\toterr]}
		+ \poldiam \sum_{\runalt=\start}^{\run} \iter[\step] \iter[\bbound]
		+ \sum_{\runalt=\start}^{\run} \iter[\step]^{2} \iter[\varrho]^{2}
		- \thres \sum_{\runalt=\start}^{\run} \probof{\beforeiter[\tilde\eventalt]}
\end{equation}
where we set
\begin{equation}
\curr[\varrho]^{2}
	= \poldiam^{2} \curr[\sdev]^{2} + \tfrac{3}{2}  (\vbound^{2} + \curr[\bbound]^{2} + \curr[\sdev]^{2}).
\end{equation}
Hence, combining \eqref{eq:prob-large1} and \eqref{eq:prob-large2} and invoking our stated assumptions for $\curr[\step]$, $\curr[\bbound]$ and $\curr[\sdev]$, we get
\begin{equation}
\sum_{\runalt=\start}^{\run} \probof{\iter[\tilde\eventalt]}
	\leq \frac{1}{\thres}
		\sum_{\runalt=\start}^{\run}
			\bracks{ \iter[\step] \iter[\bbound] \poldiam
				+ \iter[\step]^{2} \iter[\varrho]^{2}}
	\leq \frac{\Const}{\thres}
\end{equation}
for some $\Const \equiv \Const(\step,\runoff) > 0$ with $\lim_{\step\to0^{+}} \Const(\step,\runoff) = \lim_{\runoff\to\infty} \Const(\step,\runoff) = 0$ \revise{(since $\curr[\step] = \step / (\run+\runoff)^{\pexp}$ and $\pexp>0$)}.

Now, by choosing $\step$ sufficiently small (or $\runoff$ sufficiently large), we can ensure that $\Const/\thres < \conf$;
thus, given that the events $\iter[\tilde\eventalt]$ are disjoint for all $\runalt=\running$, we get
\(
%\begin{equation}
\probof[\big]{\union_{\runalt=\start}^{\run} \iter[\tilde\eventalt]}
	= \sum_{\runalt=\start}^{\run} \probof{\iter[\tilde\eventalt]}
	\leq \conf.
%\end{equation}
\)
In turn, this implies that
\(
%\begin{equation}
\probof{\curr[\eventalt]}
	= \probof[\big]{\comp{\init[\tilde\eventalt]} \cap \dotsi \cap \comp{\curr[\tilde\eventalt]}}
	\geq 1 - \conf,
%\end{equation}
\)
and our assertion follows.
\end{proof}


We conclude this appendix with the proof of our technical result on the events $\curr[\event]$ and $\curr[\eventalt]$:

\begin{proof}[Proof of \cref{lem:events}]
The first claim of the lemma is obvious.
For the second, we proceed inductively:

\begin{enumerate}[leftmargin=2.5em]
\item
For the base case $\run=\start$, we have $\init[\event] = \{\init \in \basin \} \supseteq \{\init \in \nhd \} = \samples$ (recall that $\init$ is initialized in $\nhd \subseteq \basin$).
Since $\eventalt_{0} = \samples$, our claim follows.

\item
Inductively, assume that $\prev[\eventalt] \subseteq \curr[\event]$ for some $\run\geq\start$.
To show that $\curr[\eventalt] \subseteq \next[\event]$,
%fix a realization in $\curr[\eventalt]$ so
suppose that $\iter[\toterr] \leq \thres$ for all $\runalt=\running,\run$.
Since $\curr[\eventalt] \subseteq \prev[\eventalt]$, this implies that $\curr[\event]$ also occurs, \ie $\iter\in\basin$ for all $\runalt=\running,\run$;
as such, it suffices to show that $\next\in\basin$.
To do so, given that $\iter\in\nhd\subseteq\basin$ for all $\runalt=\running\run$,
we readily obtain
%telescoping the bound \eqref{eq:template} over $\runalt = \running,\run$ gives
\begin{equation}
\afteriter[\energy]
	\leq \iter[\energy]
		+ \iter[\step] \iter[\snoise]
		+ \iter[\step] \iter[\sbias]
		+ \iter[\step]^{2} \iter[\second]^{2},
	\quad
	\text{for all $\runalt = \running\run$},
\end{equation}
and hence, after telescoping over $\runalt = \running,\run$, we get
\begin{equation}
\next[\energy]
	\leq \init[\energy]
		+ \curr[\mart]
%		+ \curr[\Psi]
		+ \curr[\submart]
	\leq \init[\energy]
		+ \sqrt{\curr[\toterr]}
%		+ \curr[\Psi]
		+ \curr[\toterr]
	\leq \thres
		+ \sqrt{\thres}
%		+ \thres
		+ \thres
	= 2\thres + \sqrt{\thres}.
\end{equation}
%by the inductive hypothesis.
We conclude that $\energy(\next) \leq 2\thres + \sqrt{\thres}$, \ie $\next\in\basin$, as required for the induction.
\end{enumerate}

For our third claim, note first that
\begin{align}
\label{eq:noise-tot-upd}
\curr[\toterr]
	&= (\prev[\mart] + \curr[\step]\curr[\snoise])^{2}
		+ \prev[\submart]
		+ \curr[\step] \curr[\sbias]
		+ \curr[\step]^{2} \curr[\second]^{2}
	\notag\\
	&= \prev[\toterr]
		+ 2 \curr[\step] \curr[\snoise] \prev[\mart]
		+ \curr[\step]^{2} \curr[\snoise]^{2}
		+ \curr[\step] \curr[\sbias]
		+ \curr[\step]^{2} \curr[\second]^{2},
\end{align}
so, after taking expectations, we get
\begin{align}
\exof{\curr[\toterr] \given \curr[\filter]}
	= \prev[\toterr]
		+ 2 \prev[\mart] \curr[\step] \exof{\curr[\snoise] \given \curr[\filter]}
		+ \exof{
			\curr[\step]^{2} \curr[\snoise]^{2}
			+ \curr[\step] \curr[\sbias]
			+ \curr[\step]^{2} \curr[\second]^{2}
			\given \curr[\filter] }
	\geq \prev[\toterr],
\end{align}
\ie $\curr[\toterr]$ is a submartingale.
To proceed, let $\curr[\tilde\toterr] = \curr[\toterr] \one_{\prev[\eventalt]}
$ so
\begin{align}
\label{eq:noise-tot-cond1}
\curr[\tilde\toterr]
	= \curr[\toterr] \one_{\prev[\eventalt]}
	&= \prev[\toterr] \one_{\prev[\eventalt]}
		+ (\curr[\toterr] - \prev[\toterr]) \one_{\prev[\eventalt]}
	\notag\\
	&= \prev[\toterr] \one_{\beforeprev[\eventalt]}
		- \prev[\toterr] \one_{\prev[\tilde\eventalt]}
		+ (\curr[\toterr] - \prev[\toterr]) \one_{\prev[\eventalt]},
	\notag\\
	&= \prev[\tilde\toterr]
		+ (\curr[\toterr] - \prev[\toterr]) \one_{\prev[\eventalt]}
		- \prev[\toterr] \one_{\prev[\tilde\eventalt]},
\end{align}
where we used the fact that $\prev[\eventalt] = \beforeprev[\eventalt] \setminus \prev[\tilde\eventalt]$ so $\one_{\prev[\eventalt]} = \one_{\beforeprev[\eventalt]} - \one_{\prev[\tilde\eventalt]}$ (since $\prev[\eventalt] \subseteq \beforeprev[\eventalt]$).
Then, \eqref{eq:noise-tot-upd} yields
\begin{align}
\curr[\toterr] - \prev[\toterr]
		= 2 \prev[\mart] \curr[\step] \curr[\snoise]
		+ \curr[\step]^{2} \curr[\snoise]^{2}
		+ \curr[\step] \curr[\sbias]
		+ \curr[\step]^{2} \curr[\second]^{2}
\end{align}
and hence, given that $\prev[\eventalt]$ is $\curr[\filter]$-measurable, we get:
\begin{subequations}
\begin{align}
\exof{(\curr[\toterr] - \prev[\toterr]) \one_{\prev[\eventalt]}}
	&\label{eq:noise-tot-zero}
		= 2 \exof{\curr[\step] \prev[\mart]\curr[\snoise] \one_{\prev[\eventalt]}}
	\\
	&\label{eq:noise-tot-noise}
		+ \exof{\curr[\step]^{2} \curr[\snoise]^{2} \one_{\prev[\eventalt]}}
	\\
	&\label{eq:noise-tot-signal}
		+ \exof{(\curr[\step]\curr[\sbias] + \curr[\step]^{2} \curr[\second]^{2}) \one_{\prev[\eventalt]}}.
\end{align}
\end{subequations}
However, since $\prev[\eventalt]$ and $\prev[\mart]$ are both $\curr[\filter]$-measurable, we have the following estimates:

\begin{enumerate}
\item
For the noise term in \eqref{eq:noise-tot-zero}, we have:
\begin{equation}
\exof{\prev[\mart] \curr[\snoise] \one_{\prev[\eventalt]}}
	= \exof{\prev[\mart] \one_{\prev[\eventalt]} \exof{\curr[\snoise] \given \curr[\filter]}}
	= 0.
\end{equation}

\item
The term \eqref{eq:noise-tot-noise} is where the reduction to $\prev[\eventalt]$ kicks in;
indeed, we have:
\begin{align}
\exof{\curr[\snoise]^{2} \one_{\prev[\eventalt]}}
	&= \exof{\one_{\prev[\eventalt]} \exof{ \abs{\braket{\curr - \eq}{\curr[\noise]}}^{2}
		\given \curr[\filter]} }
	\notag\\
	&\leq \exof{\one_{\prev[\eventalt]} \norm{\curr - \eq}^{2} \exof{\dnorm{\curr[\noise]}^{2}
		\given \curr[\filter]}}
	\explain{by Cauchy\textendash Schwarz}
	\\
%	&\leq \exof{\one_{\curr[\event]} \norm{\curr - \eq}^{2} \exof{\dnorm{\curr[\noise]}^{2}
%		\given \curr[\filter]}}
%	\explain{because $\prev[\eventalt] \subseteq \curr[\event]$}
%	\\
	&\leq \poldiam^{2} \curr[\sdev]^{2}.
%	\explain{by \cref{eq:variance}}
\end{align}

\item
Finally, for the term \eqref{eq:noise-tot-signal}, we have:
\begin{align}
\label{eq:noise-tot-term1}
\exof{\curr[\second]^{2} \one_{\prev[\eventalt]}}
	\leq \tfrac{3}{2}  \bracks{ \vbound^{2} + \curr[\bbound]^{2} + \curr[\sdev]^{2} }
\end{align}
where we used the bound $\dnorm{\gvalue(\policy)} \leq \vbound$.
Likewise, $\curr[\sbias] \one_{\prev[\eventalt]} \leq \poldiam \curr[\bbound]$, so
\begin{align}
\label{eq:noise-tot-term2}
\eqref{eq:noise-tot-signal}
	&\leq \curr[\step] \poldiam \curr[\bbound]
		+ \tfrac{3}{2}  \curr[\step]^{2} (\vbound^{2} + \curr[\bbound]^{2} + \curr[\sdev]^{2})
\end{align}
\end{enumerate}
Thus, putting together all of the above, we obtain:
\begin{equation}
\exof{(\curr[\toterr] - \prev[\toterr]) \one_{\prev[\eventalt]}}
	\leq \curr[\step] \poldiam \curr[\bbound]
		+ \curr[\step]^{2} \poldiam^{2} \curr[\sdev]^{2}
		+ \tfrac{3}{2}  \curr[\step]^{2} (\vbound^{2} + \curr[\bbound]^{2} + \curr[\sdev]^{2})
\end{equation}
Going back to \eqref{eq:noise-tot-cond1}, we have $\prev[\toterr] > \thres$ if $\prev[\tilde\eventalt]$ occurs, so the last term becomes
\begin{equation}
\label{eq:noise-tot-term3}
\exof{\prev[\toterr] \one_{\prev[\tilde\eventalt]}}
	\geq \thres \exof{\one_{\prev[\tilde\eventalt]}}
	= \thres \probof{\prev[\tilde\eventalt]}.
\end{equation}
Our claim then follows by combining \cref{eq:noise-tot-cond1,eq:noise-tot-term1,,eq:noise-tot-term2,eq:noise-tot-term3}.
\end{proof}



%----------------------------------------------------------------------
%% Subsequence
%----------------------------------------------------------------------
\subsection{Extraction of a convergent subsequence}

Our next step is to show that any realization $\curr$ of \eqref{eq:PG} that is contained in $\basin$ admits a subsequence $\state_{\run_{\runalt}}$ converging to $\eq$.
%Formally, we have:

\begin{proposition}
\label{prop:subsequence}
Let $\eq$ be a stable Nash policy, and let $\curr$ be the sequence of play generated by \eqref{eq:PG} with step-size and policy gradient estimates such that $\pexp + \bexp >1$ and $\pexp - \sexp > 1/2$ as per \eqref{eq:schedule}.
Then $\curr$ admits a subsequence $\state_{\run_{\runalt}}$ that converges to $\eq$ \acl{wp1} on the event $\event = \intersect_{\run} \curr[\event] = \{\curr\in\basin \text{ for all } \run=\running\}$.
\end{proposition}

\begin{proof}
Let $\mathcal{Q} = \{ \text{$\curr\in\basin$ for all $\run$} \} \cap \{ \liminf_{\run} \norm{\curr - \eq} > 0 \}$ denote the event that $\curr$ is contained in $\basin$ but the sequence $\curr$ does not admit a subsequence converging to $\eq$.
We will show that $\probof{\mathcal{Q}} = 0$.

Indeed, assume ad absurdum that $\probof{\mathcal{Q}} > 0$.
Hence, \acl{wp1} on $\mathcal{Q}$, there exists some positive constant $\const>0$ (again, possibly random) such that $\braket{\gvalue(\curr)}{\curr - \eq} \leq -\const < 0$ for all $\run$.
Thus, going back to \eqref{eq:energy}, we get
\begin{equation}
\next[\energy]
	\leq \curr[\energy]
		- \curr[\step] \const
		+ \curr[\step] \curr[\snoise]
		+ \curr[\step] \curr[\sbias]
		+ \curr[\step]^{2} \curr[\second]^{2},
\end{equation}
so if we let $\curr[\tau] = \sum_{\runalt=\start}^{\run} \iter[\step]$ and telescope the above, we obtain the bound
\begin{equation}
\label{eq:energy-bound1}
\next[\energy]
	\leq \init[\energy]
		- \curr[\tau] \bracks*{
			\const
			- \frac{\curr[\mart]}{\curr[\tau]}
			- \frac{\curr[\submart]}{\curr[\tau]}
%			- \frac{\sum_{\runalt=\start}^{\run} \curr[\step]\curr[\snoise]}{\curr[\tau]}
%			- \frac{\sum_{\runalt=\start}^{\run} \curr[\step]\curr[\sbias]}{\curr[\tau]}
%			- \frac{\sum_{\runalt=\start}^{\run} \curr[\step]^{2} \curr[\second]^{2}}{\curr[\tau]}
		}	% end bracks
\end{equation}
with $\curr[\snoise]$,
$\curr[\sbias]$
and
$\curr[\second]$ given by \eqref{eq:errors},
and
$\curr[\mart] = \sum_{\runalt=\start}^{\run} \iter[\step] \iter[\snoise]$,
$\curr[\submart] = \sum_{\runalt=\start}^{\run} \bracks{\iter[\step] \iter[\sbias] + \iter[\step]^{2} \iter[\second]^{2}}$
defined as in \eqref{eq:err-agg}.
Also, \eqref{eq:errorbounds} readily gives
\begin{equation}
\sum_{\run=\start}^{\infty} \exof{\curr[\step]^{2} \curr[\snoise]^{2} \given \curr[\filter]}
	\leq \sum_{\run=\start}^{\infty} \curr[\step]^{2} \exof{ \norm{\curr - \eq}^{2} \dnorm{\curr[\noise]}^{2} \given \curr[\filter]}
	\leq \poldiam^{2} \sum_{\run=\start}^{\infty} \curr[\step]^{2} \curr[\sdev]^{2}
	< \infty
\end{equation}
so, by the strong law of large numbers for martingale difference sequences \citep[Theorem 2.18]{HH80}, we conclude that $\curr[\mart] / \curr[\tau]$ converges to $0$ \acl{wp1}.
In a similar vein, for the submartingale $\curr[\submart]$ we have
\begin{align}
\exof{\curr[\submart]}
	&= \sum_{\runalt=\start}^{\run}
		\iter[\step] \iter[\sbias]
		+ \sum_{\runalt=\start}^{\run} \iter[\step]^{2} \exof{\iter[\second]^{2}}
%	\notag\\
	\leq \poldiam \sum_{\runalt=\start}^{\run} \iter[\step] \iter[\bbound]
		+ \frac{3}{2} \sum_{\runalt=\start}^{\run} \iter[\step]^{2}
			\bracks{\vbound^{2} + \iter[\bbound]^{2} + \iter[\sdev]^{2}},
\end{align}
so, by \eqref{eq:errorbounds} and the stated conditions for the method's step-size and bias/noise parameters, it follows that $\curr[\submart]$ is bounded in $L^{1}$.
Therefore, by Doob's submartingale convergence theorem \citep[Theorem~2.5]{HH80}, we further deduce that $\curr[\submart]$ converges \acl{wp1} to some (finite) random variable $\submart_{\infty}$.
%which implies in turn that $\lim_{\run\to\infty} \curr[\submart] / \curr[\tau] = 0$ \as.

Going back to \eqref{eq:energy-bound1} and letting $\run\to\infty$, the above shows that $\curr[\energy] \to -\infty$ \acl{wp1} on $\mathcal{Q}$.
Since $\energy$ is nonnegative by construction and $\probof{\mathcal{Q}} > 0$ by assumption, we obtain a contradiction and our proof is complete.
\end{proof}



%----------------------------------------------------------------------
%% Energy values
%----------------------------------------------------------------------
\subsection{Convergence of the energy values}

Our last auxiliary result concerns the convergence of the values of the dual energy function $\energy$.
We encode this as follows.

\begin{proposition}
\label{prop:energy}
If \eqref{eq:PG} is run with assumptions as in \cref{prop:contain}, there exists a finite random variable $\energy_{\infty}$ such that
\begin{equation}
\label{eq:energy-conv}
\probof*{\text{$\curr[\energy] \to \energy_{\infty}$ as $\run\to\infty$} \given \text{$\curr\in\basin$ for all $\run$}}
	= 1.
\end{equation}
\end{proposition}

\begin{proof}
Let $\curr[\event] = \braces*{\iter\in\basin \; \text{for all $\runalt=\running,\run$}}$ be defined as in \eqref{eq:events}, and let $\curr[\tilde\energy] = \one_{\curr[\event]} \curr[\energy]$.
Then, by the energy inequality \eqref{eq:template} and the fact that $\next[\event] \subseteq \curr[\event]$, we get
\begin{align}
\next[\tilde\energy]
	= \one_{\next[\event]} \next[\energy]
	&\leq \one_{\curr[\event]} \next[\energy]
	\notag\\
	&\leq \one_{\curr[\event]} \curr[\energy]
		+ \one_{\curr[\event]} \curr[\step] \braket{\gvalue(\curr)}{\curr - \eq}
		+ \parens[\big]{
			\curr[\step] \curr[\snoise]
			+ \curr[\step] \curr[\sbias]
			+ \curr[\step]^{2} \curr[\second]^{2}
			} % end bracks
			\one_{\curr[\event]}
	\notag\\
	&\leq \curr[\tilde\energy]
		+ \curr[\step] \one_{\curr[\event]} \curr[\snoise]
		+ \parens[\big]{
			\curr[\step] \curr[\sbias]
			+ \curr[\step]^{2} \curr[\second]^{2}
			} % end bracks
			\one_{\curr[\event]},
\end{align}
where we used the fact that that $\braket{\gvalue(\iter)}{\iter - \eq} \leq 0$ for all $\runalt = \running,\run$ if $\curr[\event]$ occurs.
Since $\curr[\event]$ is $\curr[\filter]$-measurable, conditioning on $\curr[\filter]$ and taking expectations yields
\begin{align}
\exof{\next[\tilde\energy] \given \curr[\filter]}
	&\leq \curr[\tilde\energy]
		+ \curr[\step] \one_{\curr[\event]} \exof{\curr[\snoise] \given \curr[\filter]}
		+ \one_{\curr[\event]} \curr[\step] \curr[\sbias]
		+ \one_{\curr[\event]} \exof{\curr[\step]^{2} \curr[\second]^{2} \given \curr[\filter]}
	\notag\\
	&\leq \curr[\tilde\energy]
		+ \curr[\step] \poldiam \curr[\bbound]
		+ \curr[\step] \curr[\sbias]
		+ \exof{\curr[\step]^{2} \curr[\second]^{2} \given \curr[\filter]}
	\notag\\
	&\leq \curr[\tilde\energy]
		+ \curr[\step] \poldiam \curr[\bbound]
		+ \frac{3}{2} \bracks[\big]{ \vbound^{2} + \curr[\bbound]^{2} + \curr[\sdev]^{2}}.
\end{align}
By our step-size assumptions, we have $\sum_{\run} \curr[\step]^{2} (1 + \curr[\bbound]^{2} + \curr[\sdev]^{2}) < \infty$ and $\sum_{\run} \curr[\step] \curr[\bbound] < \infty$, which means that $\curr[\tilde\energy]$ is an almost supermartingale with almost surely summable increments, \ie
\begin{equation}
\sum_{\run=\start}^{\infty} \bracks*{\exof{\next[\tilde\energy] \given \curr[\filter]} - \curr[\tilde\energy]}
	< \infty
	\quad
	\text{\acl{wp1}}
\end{equation}
Therefore, by Gladyshev's lemma \citep[p.~49]{Pol87}, we conclude that $\curr[\tilde\energy]$ converges almost surely to some (finite) random variable $\energy_{\infty}$.
Since $\one_{\curr[\event]} = 1$ for all $\run$ if and only if $\curr\in\basin$ for all $\run$, we conclude that $\probof*{\text{$\curr[\energy]$ converges} \given \text{$\curr\in\basin$ for all $\run$}} = \probof{\text{$\curr[\tilde\energy]$ converges}}= 1$, and our claim follows.
% \KL{Maybe $\probof{\text{$\curr[\tilde\energy]$ converges} \given  \text{$\curr\in\basin$ for all $\run$}}$}
% \PM{Recall the definition of $\curr[\tilde\energy]$, it subsumes this.}
\end{proof}


%----------------------------------------------------------------------
%% Proof of theorem
%----------------------------------------------------------------------
\subsection{Putting everything together}
We are now in a position to prove \cref{thm:stable,cor:stable}.


\begin{proof}[Proof of \cref{thm:stable}]
Let $\event = \intersect_{\run} \curr[\event] = \{\text{$\curr\in\basin$ for all $\run$} \}$ denote the event that $\curr$ lies in $\basin$ for all $\run$.
By \cref{prop:contain}, if $\init$ is initialized within the neighborhood $\nhd$ defined in \eqref{eq:nhd-init}, we have $\probof{\event \given \init\in\nhd} \geq 1-\thres$, noting also that the neighborhood $\nhd$ is independent of the required confidence level $\thres$.
Then, by \cref{prop:subsequence,prop:energy}, it follows that
\begin{enumerate*}
[\itshape a\upshape)]
\item
$\liminf_{\run} \norm{\curr - \eq} = 0$;
and
\item
$\curr[\energy]$ converges,
\end{enumerate*}
both events occurring \acl{wp1} on the set $\event \cap \{ \init\in\nhd \}$.
We thus conclude that $\lim_{\run\to\infty} \curr[\energy] = 0$ and hence
\begin{align*}
\probof{ \curr\to\eq \given \init \in \nhd}
	&\geq \probof{ \event \cap \{ \curr\to\eq \} \given \init \in \nhd}
	\notag\\
	&= \probof{ \curr\to\eq \given \init \in \nhd, \event}
		\times \probof{\event \given \init \in \nhd}
	\geq 1 - \conf,
\end{align*}
and our proof is complete.
\end{proof}


\begin{proof}[Proof of \cref{cor:stable}]
For \cref{mod:full,mod:stoch}, taking $\bexp = \infty, \sexp = 0$, we obtain $\pexp > 1/2$. Since we have that $\sum_{\run=1}^{\infty}\step_\run = \infty$, we get that $\pexp \leq 1$, i.e., $\pexp \in (1/2,1]$.

For \cref{mod:value}, we have that $\curr[\bbound] = \bigoh(\curr[\expar])$ and $\curr[\sdev] = \bigoh(1/\sqrt{\curr[\expar]})$, i.e., $\bexp = \rexp$ and $\sexp = \rexp/2$. Now, since $\pexp \leq 1$, $\pexp + \bexp >1$ and $\pexp - \sexp > 1/2$, we obtain that $\pexp \in (2/3,1]$ and $(1-p)/2 < r/2 < p - 1/2$.
\end{proof}